\documentclass[10pt,a4paper]{amsart}
\usepackage{amsmath,amssymb,amsthm,graphicx}
\usepackage{verbatim}
\newenvironment{reference}{\comment}{\endcomment}
%\newenvironment{reference}{}{}
\newenvironment{slogan}{\comment}{\endcomment}
\newenvironment{history}{\comment}{\endcomment}

% For commutative diagrams we use Xy-pic
\usepackage[all]{xy}

% We use 2cell for 2-commutative diagrams.
\xyoption{2cell}
% Theorem environments.
%
\theoremstyle{plain}
\newtheorem{theorem}[subsection]{Theorem}
\newtheorem{proposition}[subsection]{Proposition}
\newtheorem{lemma}[subsection]{Lemma}

\theoremstyle{definition}
\newtheorem{definition}[subsection]{Definition}
\newtheorem{example}[subsection]{Example}
\newtheorem{exercise}[subsection]{Exercise}
\newtheorem{situation}[subsection]{Situation}

\theoremstyle{remark}
\newtheorem{remark}[subsection]{Remark}
\newtheorem{remarks}[subsection]{Remarks}

\numberwithin{equation}{subsection}

% Macros
%
\def\lim{\mathop{\mathrm{lim}}\nolimits}
\def\colim{\mathop{\mathrm{colim}}\nolimits}
\def\Spec{\mathop{\mathrm{Spec}}}
\def\Hom{\mathop{\mathrm{Hom}}\nolimits}
\def\Ext{\mathop{\mathrm{Ext}}\nolimits}
\def\SheafHom{\mathop{\mathcal{H}\!\mathit{om}}\nolimits}
\def\SheafExt{\mathop{\mathcal{E}\!\mathit{xt}}\nolimits}
\def\Sch{\mathit{Sch}}
\def\Mor{\mathop{\mathrm{Mor}}\nolimits}
\def\Ob{\mathop{\mathrm{Ob}}\nolimits}
\def\Sh{\mathop{\mathit{Sh}}\nolimits}
\def\NL{\mathop{N\!L}\nolimits}
\def\CH{\mathop{\mathrm{CH}}\nolimits}
\def\proetale{{pro\text{-}\acute{e}tale}}
\def\etale{{\acute{e}tale}}
\def\QCoh{\mathit{QCoh}}
\def\Ker{\mathop{\mathrm{Ker}}}
\def\Im{\mathop{\mathrm{Im}}}
\def\Coker{\mathop{\mathrm{Coker}}}
\def\Coim{\mathop{\mathrm{Coim}}}

% Boxtimes
%
\DeclareMathSymbol{\boxtimes}{\mathbin}{AMSa}{"02}

%
% Macros for moduli stacks/spaces
%
\def\QCohstack{\mathcal{QC}\!\mathit{oh}}
\def\Cohstack{\mathcal{C}\!\mathit{oh}}
\def\Spacesstack{\mathcal{S}\!\mathit{paces}}
\def\Quotfunctor{\mathrm{Quot}}
\def\Hilbfunctor{\mathrm{Hilb}}
\def\Curvesstack{\mathcal{C}\!\mathit{urves}}
\def\Polarizedstack{\mathcal{P}\!\mathit{olarized}}
\def\Complexesstack{\mathcal{C}\!\mathit{omplexes}}
% \Pic is the operator that assigns to X its picard group, usage \Pic(X)
% \Picardstack_{X/B} denotes the Picard stack of X over B
% \Picardfunctor_{X/B} denotes the Picard functor of X over B
\def\Pic{\mathop{\mathrm{Pic}}\nolimits}
\def\Picardstack{\mathcal{P}\!\mathit{ic}}
\def\Picardfunctor{\mathrm{Pic}}
\def\Deformationcategory{\mathcal{D}\!\mathit{ef}}
\usepackage{cleveref}
\crefname{originalreference}{original reference}{original references}
\Crefname{originalreference}{Original reference}{Original references}
\makeatletter
\newlabel{sites-lemma-points-above-point}{{04H3}{source}{}{sourcepointer-1}{}}
\newlabel{sites-lemma-points-above-point@cref}{{[originalreference][0][]04H3}{[1][0][]source}{}{}{}}
\makeatother
\crefname{originallemma}{original lemma}{original lemmas}
\Crefname{originallemma}{Original lemma}{Original lemmas}
\makeatletter
\newlabel{lemma-extension-by-zero}{{03DI}{source}{}{sourcepointer-2}{}}
\newlabel{lemma-extension-by-zero@cref}{{[originallemma][0][]03DI}{[1][0][]source}{}{}{}}
\makeatother
\crefname{originallemma}{original lemma}{original lemmas}
\Crefname{originallemma}{Original lemma}{Original lemmas}
\makeatletter
\newlabel{lemma-stalk-exact-modules}{{04EQ}{source}{}{sourcepointer-3}{}}
\newlabel{lemma-stalk-exact-modules@cref}{{[originallemma][0][]04EQ}{[1][0][]source}{}{}{}}
\makeatother
\crefname{originalreference}{original reference}{original references}
\Crefname{originalreference}{Original reference}{Original references}
\makeatletter
\newlabel{sites-section-points}{{00Y3}{source}{}{sourcepointer-4}{}}
\newlabel{sites-section-points@cref}{{[originalreference][0][]00Y3}{[1][0][]source}{}{}{}}
\makeatother
\title{Stalks of extension by zero: \$(j\_\{U!\}\textbackslash{}mathcal\{G\})\_p = \textbackslash{}bigoplus\_q \textbackslash{}mathcal\{G\}\_q\$ over the points \$q\$ of \$\textbackslash{}mathcal\{C\}/U\$ lying over \$p\$}
\author{The Stacks Project Authors}
\date{stacks-project:sites-modules.tex:local-snapshot-20261012; GNU FDL 1.2 or later}
\begin{document}
\maketitle
\noindent\textit{Editorial note.} All delivered mathematical text is the author's own, copied verbatim from the pinned Stacks Project chapter file sites-modules.tex; no statement or proof step was rewritten, reordered or removed.\par\medskip
\noindent\textit{Editorial note.} Span recomputation by label pairing (LANE\_BRIEF 10.2): in this chapter the author's proof environment follows the statement environment as a standalone block; the spec covers the complete environment pair (statement span, proof span) as recorded in the pool candidate and re-verified against the pinned source (sm-twins-20261014.json span sanity: all OK).\par\medskip
\noindent\textit{Editorial note.} Occupation scan: combined claim\_key + statement\_anchor + generic-line scan over all registration files (active and withdrawn) found hard occupation only from the 47 pre-existing lane-s11 S11-sites-mod-* cards (same pinned snapshot, delivered) and cross-chapter citation noise, 0 other-lane hits for sites-modules.tex (sm-occupied.json, 2026-10-14); the RESUME-pool dispatch unlocks the remaining face.\par\medskip
\ifdefined\hypertarget\hypertarget{sourcepointer-1}{}\fi
\noindent\textit{Source reference.} Original reference 04H3: sources/stacks-chow/tags.txt, source lines 5854--5854. Cross-chapter standard prerequisite cited by the author inside the counted span; it is used as a pointer and its own proof is not part of this item.\par\medskip
\ifdefined\hypertarget\hypertarget{sourcepointer-2}{}\fi
\noindent\textit{Source reference.} Original lemma 03DI: sources/stacks-chow/tags.txt, source lines 4421--4421. Cross-chapter standard prerequisite cited by the author inside the counted span; it is used as a pointer and its own proof is not part of this item.\par\medskip
\ifdefined\hypertarget\hypertarget{sourcepointer-3}{}\fi
\noindent\textit{Source reference.} Original lemma 04EQ: sources/stacks-chow/tags.txt, source lines 5771--5771. Cross-chapter standard prerequisite cited by the author inside the counted span; it is used as a pointer and its own proof is not part of this item.\par\medskip
\ifdefined\hypertarget\hypertarget{sourcepointer-4}{}\fi
\noindent\textit{Source reference.} Original reference 00Y3: sources/stacks-chow/tags.txt, source lines 1241--1241. Cross-chapter standard prerequisite cited by the author inside the counted span; it is used as a pointer and its own proof is not part of this item.\par\medskip
\begin{lemma}
\label{lemma-stalk-j-shriek}
Let $(\mathcal{C}, \mathcal{O})$ be a ringed site.
Let $p$ be a point of $\mathcal{C}$. Let $U$ be an object of $\mathcal{C}$.
For $\mathcal{G}$ in $\textit{Mod}(\mathcal{O}_U)$ we have
$$
(j_{U!}\mathcal{G})_p =
\bigoplus\nolimits_q \mathcal{G}_q
$$
where the coproduct is over the points $q$ of $\mathcal{C}/U$
lying over $p$, see
Sites, Lemma \ref{sites-lemma-points-above-point}.
\end{lemma}

\begin{proof}
We use the description of $j_{U!}\mathcal{G}$ as the sheaf associated
to the presheaf
$V \mapsto
\bigoplus\nolimits_{\varphi \in \Mor_\mathcal{C}(V, U)}
\mathcal{G}(V/_\varphi U)$
of
Lemma \ref{lemma-extension-by-zero}.
The stalk of $j_{U!}\mathcal{G}$ at $p$ is equal to the
stalk of this presheaf, see
Lemma \ref{lemma-stalk-exact-modules}.
Let $u : \mathcal{C} \to \textit{Sets}$ be the functor corresponding
to $p$ (see Sites, Section \ref{sites-section-points}).
Hence we see that
$$
(j_{U!}\mathcal{G})_p = \colim_{(V, y)}
\bigoplus\nolimits_{\varphi : V \to U} \mathcal{G}(V/_\varphi U)
$$
where the colimit is taken in the category of abelian groups.
To a quadruple $(V, y, \varphi, s)$ occurring in this colimit, we can assign
$x = u(\varphi)(y) \in u(U)$. Hence we obtain
$$
(j_{U!}\mathcal{G})_p =
\bigoplus\nolimits_{x \in u(U)}
\colim_{(\varphi : V \to U, y), \ u(\varphi)(y) = x} \mathcal{G}(V/_\varphi U).
$$
This is equal to the expression of the lemma by the description
of the points $q$ lying over $x$ in
Sites, Lemma \ref{sites-lemma-points-above-point}.
\end{proof}

\end{document}
